\documentclass[11pt,a4paper]{article}
\usepackage{fontspec,amsmath,amssymb,amsthm,microtype}
\setmainfont{lmroman10-regular.otf}[BoldFont=lmroman10-bold.otf,ItalicFont=lmroman10-italic.otf,BoldItalicFont=lmroman10-bolditalic.otf]
\usepackage{xcolor}
\usepackage[most]{tcolorbox}
\usepackage{mathrsfs}
\providecommand{\lt}{<}
\providecommand{\gt}{>}
\setmonofont{DejaVuSansMono.ttf}[BoldFont=DejaVuSansMono-Bold.ttf,ItalicFont=DejaVuSansMono-Oblique.ttf,BoldItalicFont=DejaVuSansMono-BoldOblique.ttf,Scale=0.86]
\newfontfamily\leanfallback{STIXTwoMath-Regular.otf}[Scale=0.86]
\definecolor{leanrule}{gray}{0.62}
\definecolor{leanback}{gray}{0.962}
\definecolor{leantitle}{gray}{0.22}
\definecolor{leanmuted}{gray}{0.38}
\newtcolorbox{leanbox}[1]{enhanced,breakable,colback=leanback,colframe=leanrule,
  boxrule=0.4pt,arc=0.6mm,left=2.4mm,right=2.4mm,top=1.8mm,bottom=1.8mm,
  before skip=0.8em,after skip=1em,fonttitle=\small\bfseries,coltitle=leantitle,
  colbacktitle=leanback,titlerule=0.3pt,title={#1}}
\newlength{\leancharwidth}
\newenvironment{leancode}{\par\smallskip\begingroup\ttfamily\footnotesize\raggedright
  \settowidth{\leancharwidth}{M}\setlength{\parskip}{0pt}\setlength{\parindent}{0pt}
  \frenchspacing}{\par\endgroup\smallskip}
\newcommand{\leanline}[2]{\hangindent=\dimexpr#1\leancharwidth+4\leancharwidth\relax
  \hangafter=1\noindent\hspace*{#1\leancharwidth}#2\par}
\newcommand{\leansource}[1]{\par{\footnotesize\color{leanmuted}#1}\par}
\newcommand{\leanlabel}[1]{\par\smallskip\noindent{\small\bfseries #1}\par}
\usepackage[margin=25mm]{geometry}
\usepackage[unicode,hidelinks]{hyperref}
\hypersetup{pdfauthor={LangevinLean Project}}
\setlength{\emergencystretch}{3em}
\setlength{\parskip}{0.4em}
\setlength{\parindent}{0pt}
\allowdisplaybreaks[2]
\raggedbottom
\title{\texorpdfstring{Accumulation of quadratic local errors under OBABO contraction}{Accumulation of quadratic local errors under OBABO contraction}}
\author{LangevinLean Project}
\date{}
\begin{document}\maketitle
Draft mathematical exposition; these notes may contain mistakes.


The contracting metric factor is \(r=\sqrt{1-mh/(3\gamma)}\). These scalar lemmas prove its strict contraction and a lower bound on its gap, then simplify the proved local-error bound into an explicit first-order bias coefficient.



The contraction gap is of first order in the step, while local kernel error is of second order. Their ratio produces the explicit linear-step global bias, with all friction and dimension factors visible.


\section{\texorpdfstring{Quantify the contraction gap}{Quantify the contraction gap}}


Let \(m,\gamma,h>0\) and put \(s=mh/(3\gamma)\). Then \(s>0\), so \(\sqrt{1-s}<1\) whenever the square root has its nonnegative argument interpretation. In the admitted step range, \(s\le1\); write \(r=\sqrt{1-s}\), so \(0\le r<1\) and \(r^2=1-s\).



The identity \((1-r)(1+r)=s\), or equivalently the nonnegative square \((1-r)^2\ge0\), gives

\[1-r=\frac{s}{1+r}\ge\frac s2=\frac{mh}{6\gamma}.\]

Here \(1+r\le2\) and all denominators are positive. The strict-contraction statement requires positive \(m,\gamma,h\); the displayed gap bound also requires the stated nonnegative square-root argument.

\begin{leanbox}{Lean: \texttt{obaboContractionRate\_\allowbreak{}lt\_\allowbreak{}one}, \texttt{obabo\_\allowbreak{}rate\_\allowbreak{}gap\_\allowbreak{}bound}}
\leansource{\texttt{obaboContractionRate\_\allowbreak{}lt\_\allowbreak{}one} \quad KineticLangevin/\allowbreak{}OBABOBiasScalar.lean, lines 12–17}
\begin{leancode}
\leanline{0}{theorem\ obaboContractionRate\_lt\_one\ (m\ γ\ h\ :\ ℝ)\ (hm\ :\ 0\ <\ m)\ (hγ\ :\ 0\ <\ γ)\ (hh\ :\ 0\ <\ h)\ :}
\leanline{4}{(obaboContractionRate\ m\ γ\ h\ :\ ℝ)\ <\ 1}
\end{leancode}
\leansource{\texttt{obabo\_\allowbreak{}rate\_\allowbreak{}gap\_\allowbreak{}bound} \quad KineticLangevin/\allowbreak{}OBABOBiasScalar.lean, lines 19–29}
\begin{leancode}
\leanline{0}{theorem\ obabo\_rate\_gap\_bound\ (m\ γ\ h\ :\ ℝ)\ (hm\ :\ 0\ <\ m)\ (hγ\ :\ 0\ <\ γ)\ (hh\ :\ 0\ <\ h)}
\leanline{4}{(hq\ :\ 0\ ≤\ 1\ -\ m\ *\ h\ /\ (3\ *\ γ))\ :}
\leanline{4}{m\ *\ h\ /\ (6\ *\ γ)\ ≤\ 1\ -\ (obaboContractionRate\ m\ γ\ h\ :\ ℝ)}
\end{leancode}
\leanlabel{How the Lean proof follows the argument}
\texttt{obabo\allowbreak{}Contraction\allowbreak{}Rate\_\allowbreak{}lt\_\allowbreak{}one} gives the formal counterpart of this step. The main intermediate facts are \texttt{hp} in that order. The final step is the algebraic normalization or ordered arithmetic displayed in the proof.




\texttt{obabo\_\allowbreak{}rate\_\allowbreak{}gap\_\allowbreak{}bound} gives the formal counterpart of this step. The main intermediate facts are \texttt{hr}, \texttt{hs}, \texttt{hprod}, \texttt{hle} in that order. Its remaining equalities are discharged by the simplifications displayed in the proof.
\end{leanbox}

\section{\texorpdfstring{Simplify the local-error bound and divide by the proved gap}{Simplify the local-error bound and divide by the proved gap}}


The nonnegative square-root product identity yields \(b_h=h^2\sqrt d(4\sqrt M+Mh/2)\). Dividing the local-error bound by the proved lower bound on \(1-r\) gives

\[\frac{b_h}{1-r}\le\frac{6\gamma}{m}h\sqrt d(4\sqrt M+Mh/2).\]

Thus the accumulated bias is \(O(h\sqrt d)\). Restricting \(h\) to its stated upper bound makes the remaining coefficient depend only on \(m,M,\gamma\). This is arithmetic applied to an independently proved local map error and contraction rate.



The local cost is proportional to \(h^2\sqrt d\), while the contraction gap is bounded below by a constant times \(mh/\gamma\). Dividing cancels one power of \(h\). The remaining terms give the stated linear-step bias with all physical constants retained.

\begin{leanbox}{Lean: \texttt{obaboLocalBudget\_\allowbreak{}simplify}, \texttt{obaboLocalBudget\_\allowbreak{}div\_\allowbreak{}gap}}
\leansource{\texttt{obaboLocalBudget\_\allowbreak{}simplify} \quad KineticLangevin/\allowbreak{}OBABOBiasScalar.lean, lines 31–39}
\begin{leancode}
\leanline{0}{theorem\ obaboLocalBudget\_simplify\ (M\ h\ d\ :\ ℝ)\ (hM\ :\ 0\ ≤\ M)\ :}
\leanline{4}{obaboLocalBudget\ M\ h\ d\ =\ h\ \textasciicircum{}\ 2\ *\ Real.sqrt\ d\ *\ (4\ *\ Real.sqrt\ M\ +\ M\ *\ h\ /\ 2)}
\end{leancode}
\leansource{\texttt{obaboLocalBudget\_\allowbreak{}div\_\allowbreak{}gap} \quad KineticLangevin/\allowbreak{}OBABOBiasScalar.lean, lines 41–51}
\begin{leancode}
\leanline{0}{theorem\ obaboLocalBudget\_div\_gap\ (m\ M\ γ\ h\ :\ ℝ)\ (hm\ :\ 0\ <\ m)\ (hM\ :\ 0\ ≤\ M)}
\leanline{4}{(hγ\ :\ 0\ <\ γ)\ (hh\ :\ 0\ <\ h)\ (hq\ :\ 0\ ≤\ 1\ -\ m\ *\ h\ /\ (3\ *\ γ))\ (d\ :\ ℝ)\ :}
\leanline{4}{obaboLocalBudget\ M\ h\ d\ /\ (1\ -\ (obaboContractionRate\ m\ γ\ h\ :\ ℝ))\ ≤}
\leanline{6}{(6\ *\ γ\ /\ m)\ *\ h\ *\ Real.sqrt\ d\ *\ (4\ *\ Real.sqrt\ M\ +\ M\ *\ h\ /\ 2)}
\end{leancode}
\leanlabel{How the Lean proof follows the argument}
\texttt{obabo\allowbreak{}Local\allowbreak{}Budget\_\allowbreak{}simplify} gives the formal counterpart of this step. The main intermediate facts are \texttt{hs}, \texttt{hmul} in that order. The final step is the algebraic normalization or ordered arithmetic displayed in the proof.




\texttt{obabo\allowbreak{}Local\allowbreak{}Budget\_\allowbreak{}div\_\allowbreak{}gap} gives the formal counterpart of this step. The main intermediate facts are \texttt{hgap}, \texttt{hgap\allowbreak{}Pos}, \texttt{hb}, \texttt{hdiv} in that order. The final step is the algebraic normalization or ordered arithmetic displayed in the proof.
\end{leanbox}

\end{document}
